在零槽证明中取 \(w=w_o=0\)，并令
\[
 \ell=0,\qquad g=J_++\sigma.
\]
正范数 \(\mathcal N_C\) 至多是长度为 \(K+g\) 的光滑球上的相应
范数。此情形下行之间不相乘。

对于正槽，从 \(w=w_o=0\) 开始。令
\(\ell_*=\sigma/3\)，选择槽网格 \(\eta<\sigma/6\)，
并取固定池
\[
 \mathcal P=\{p: Z^{\ell_*}/2<q_p\le Z^{\ell_*},\
                    p\notin\mathcal S\}.
\]
固定域中的素理想定理，等价地其固定
射线类形式 \cite[定理 1.1]{TZ}，给出
\(|\mathcal P|=Z^{\ell_*+o(1)}\)。因为活跃槽长度
至多为 \(\eta\) 且其相对环形支集固定，
单个不等式
\(Z^{\ell_*-\eta}>2\max_i b_i\) 使该池与
每个活跃槽窗口不相交。其指数间隔至少为 \(\sigma/6\)，
与 \(N\) 无关；只有阈值依赖于固定
窗口。对于行 \(h\ne0\)，略去整除
\(h\)、\(s\) 或任何冻结支集的池素数。每个这样的整数或理想
具有多项式范数，因此略去 \(O(\log Z)\) 个素数。
剩余集合 \(\mathcal P_h\) 的大小与
\(|\mathcal P|\) 相当，且关于该行与冻结标签一致。

为比较行 \(h\) 与 \(hp^6\)，将完整模写作 \(p^iu\)，其中
\((p,u)=1\)。CRT 与互反律给出精确公式
\[
 G(p^iu,h)
 =\mathcal R(p,u)^i\chi_p(u)^{2i}G(p^i,h)G(u,h).
\]
又有 \(G(u,hp^6)=G(u,h)\)，因为 \(p\) 是模 \(u\) 的单位
且 \(\chi_u(p^6)=1\)。
对于 \(p\nmid h\)，方程~\eqref{eq:gauss-local} 表明
将 \(h\) 替换为 \(hp^6\) 只改变 \(i=1,6,7\)。
出现的因子 \(G(p^i,h)\) 是行与
\(p\) 中的标量；所显示的剩余因子是整 \(u\) 上的一个固定射线特征
乘以一个剩余特征。因为 \(p\)
在每个活跃槽窗口之外，\(p^i\) 只分配给两个
素朴变量。剩余乘积在 \(p\) 处被穿孔，且
同一分配在两个矩形中更新 \(D_{\mathbf b}\)。
因为 \(p\nmid s\)，条件 \(s\mid u\) 仍然成立。

以下是精确的平均论证。将 \(\mathcal N_C\) 中的
正规化多项式写作 \(\mathcal H(h)\)。对每个
合格 \(p\)，局部 Gauss 求值给出
\(\mathcal H(h)=\mathcal H(hp^6)\) 加上具有
\(p\) 进列赋值 \(1,6,7\) 的抽取项。将这些赋值中每一个
分配到两个素朴变量只有有限多种方式。
Jensen 不等式——先对 \(p\) 再对此
固定有限和——给出
\[
 |\mathcal H(h)|^2\ll
 \frac1{|\mathcal P_h|}\sum_{p\in\mathcal P_h}
 \left\{|\mathcal H(hp^6)|^2+
        \sum_{i=1,6,7}\sum_{\rm allocations}
        |c_{p,i}(h)|^2|\mathcal H_{p,i}(h)|^2\right\}.
\]
\(\mathcal H_{p,i}\) 在其缩短的列
尺度上被正规化。刚才描述的局部恒等式与分配
应用于整个中心化系数，因此每个误差保持
方程~\eqref{eq:centered-coefficient-form}。

在原行权重的支集上，新行 \(hp^6\)
具有范数 \(O(Z^{K+6\ell_*})\)。固定输出行作为
\(hp^6\)（其中 \(p\in\mathcal P_h\)）只有有限多种表示：
每个这样的 \(p\) 整除该行，且所有这样的素数具有范数
至少 \(Z^{\ell_*}/2\)，而输出范数具有有界
对数长度。因此对 \(h\) 求和第一项
给出因子 \(|\mathcal P|^{-1}=Z^{-\ell_*+o(1)}\)
乘以新行上的正范数。对此主项令
\[
 \ell=\ell_*,\qquad g=J_++2\sigma.
\]
由于 \(6\ell_*=2\sigma\)，长度为 \(K+g\) 的光滑球
包含所有新行。

对于抽取项令 \(P=q_p\)。当 \(p\nmid h\) 时，
方程~\eqref{eq:gauss-local} 给出以下中心
系数。在赋值一时，旧 Gauss 和具有模一
而新的为零，故抽取给出平方系数
\(P^{-1}\)。在赋值六时，新 Gauss 和为
\(P^3(1-P^{-1})\)；其中心因子为 \(P^{-3}\)，给出
\((1-P^{-1})^2\)。在赋值七时，新 Gauss 和具有
模 \(P^3\)，中心因子 \(P^{-7/2}\) 给出
\(P^{-1}\)。行相位为 \(\overline{\chi_p(h)}\)，对
赋值一与七而言；而对赋值六为一。它
乘以两个矩形。

令 \(\ell_p=\log_ZP\)。对充分大的 \(Z\)，
\(\ell_*/2\le\ell_p\le\ell_*\)。三项的
移除长度与新移动根式长度为
\begin{equation}
 (w,w_o)=(i\ell_p,e_i\ell_p),\qquad
 i=1,6,7,\quad e_1=e_7=1,\quad e_6=0.
 \label{old-eq:2.9}
\end{equation}
两个平方因子 \(P^{-1}\) 恰好给出 \(Z^{-w_o}\)。
对赋值六项，精确因子
\(\chi_p(u)^{12}=1_{(u,p)=1}\) 由在 \(p\) 处
已经是公共的固定穿孔表示，且其固定射线相位 \(\mathcal R(p,u)^6\) 被保留。
在赋值一与七时，所显示的局部因子仍保留，且其
素数计入移动根式。在每个
误差中，冻结 \(p\)，保留其公共
列穿孔，并且仅在取正范数之后才丢弃其行资格限制。
平均值仍为
\(|\mathcal P|^{-1}\sum_{p\in\mathcal P}\)，因此关于
冻结 \(p\) 一致的界不引入素数计数因子。对每个误差
令
\[
 \ell=0,\qquad g=J_++\sigma
\]
并将其原始行直接放大为长度为
\(K+g\) 的光滑球。误差不再被放大。三种类型的
范数均满足
\begin{equation}
 0\le w_o\le w\le7\sigma/3.
 \label{old-eq:2.10}
\end{equation}

行零仅在最终光滑球处加入。令
\(Y_{\rm col}=e^{H_N}Z^{a_0}=Z^{a_0+\theta_N}\)。非空列
壳满足 \(Y_{\rm col}\ge1\)，从而 \(a_0\ge-\theta_N\)。
由方程~\eqref{eq:gauss-local}，除非 \(u\) 为
六次幂，否则 \(G(u,0)=0\)，且 \(|G(u,0)|\le q_u^{1/2}\le Y_{\rm col}^{1/2}\)。
该支集上有 \(O(Y_{\rm col}^{1/6})\) 个六次幂。
若 \(s\mid u\)，则 \(q_s^6\le q_u\le Y_{\rm col}\)，故
\(s_0\le(a_0+\theta_N)/6\)。计入所分配的除子有界
卷积损失，中心正规化后的平方贡献为
\[
 \ll Z^{-a_0}\bigl(Y_{\rm col}^{1/6}Y_{\rm col}^{1/2}\bigr)^2
          Z^{\epsilon_1}
 =Z^{a_0/3+4\theta_N/3+\epsilon_1}
 \le Z^{a_0-s_0+2\theta_N+\epsilon_1}.
\]
最后一个不等式由 \(a_0\ge-\theta_N\) 与所显示的
\(s_0\) 界得出。因此方程~\eqref{old-eq:2.7} 界定此新增
行，其中数值修正 \(2\theta_N\) 已包含在
\(C_*\xi\) 阶段余量中。这里 \(C_*\ge1\) 表示
一个阶段中累积支集误差的公共
常数，选取时与槽数 \(N\) 无关。下面的估计确立：一个这样的
选择覆盖一个阶段中的所有运算。零从不乘以
池素数。

\subsection{第二变换与较小宽度的乘积}\label{sec:plain-second-transform}

在展开平方之前恢复所有活跃槽求和。令
\(B(u)\) 表示前述某个 Gauss 多项式中完整分配的卷积系数，
包括其固定掩码
与 \(s\) 整除条件。当使用中心化时它包含整个差。
其模满足
\(q_u/Z^{a_0}\in\exp([-H_N,H_N])\)。令 \(\Phi_2\) 为
定义最终行球的非负光滑径向函数。除
素密度中的 \(Z^{o(1)}\) 损失外，待界定的范数为
\[
 Z^{-\ell}\sum_h\Phi_2(q_h/Z^{K+g})
 \left|Z^{-a_0/2}\sum_u B(u)\tau_1(u)q_u^{it}G(u,h)\right|^2.
\]
引理~\ref{lem:full-correlation} 中的 Fourier 恒等式给出其
精确展开
\[
 \begin{split}
 Z^{K+g-\ell-a_0}\sum_{u,v}
 \frac{B(u)\overline{B(v)}\tau_1(u)\overline{\tau_1(v)}
       q_u^{it}q_v^{-it}}{\sqrt{q_uq_v}}
 \sum_j F(u,v;j)
 \widehat\Phi_2\left(\frac{Z^{K+g}q_j}{q_uq_v}\right).
 \end{split}
\]
此处出现的核为
\begin{equation}
 \widehat\Phi_2\left(\frac{Z^{K+g}q_j}{q_uq_v}\right).
 \label{old-eq:2.18b}
\end{equation}
所显示的逆平方根与此核都作为整个 \(u,v\) 的
函数保留，直到下面进行整积分离。特别地
该公式包含每个活跃素数及每个
共享素数的重数。

在 \(j=0\) 时，引理~\ref{lem:full-correlation} 只留下 \(u=v\)，
且 \(F(u,u;0)=\varphi(u)\le q_u\)。受支撑的
可被 \(s\) 整除的模，在非空壳上，至多为
\[
 C\frac{Y_{\rm col}}{q_s}=C Z^{a_0-s_0+\theta_N},
\]
因为写作 \(u=sv\) 在该壳上给出 \(Y_{\rm col}/q_s\ge1\)。
除子有界系数使用其单独的 \(\epsilon_1\)
份额。因此对角指数为
\(K+g-\ell-s_0+\theta_N\)，包括 \(a_0-s_0\) 略微
为负时。从方程~\eqref{old-eq:2.7}
的余量中减去其名义部分（不含 \(\epsilon_G\)），给出
\begin{equation}
 \begin{split}
 &(K+g-\ell-s_0)
 -(a_0+\widetilde q+w_o+w+\mathcal B_c-s_0)\\
 &\quad=(A-M)-d-(K_0-K)-w_o-\mathcal B_c+g-\ell\\
 &\quad\le A-M+5\sigma/3+\delta_{{\rm fr},1}.
 \end{split}
 \label{old-eq:2.12}
\end{equation}
对于该不等式，令 \(D_0=d+K_0-K\)。贡献
\(-D_0-w_o+J_+\) 至多为 \(\delta_{{\rm fr},1}\)：当 \(J\ge0\) 时
至多为 \(-c-2w\)，而当 \(J<0\) 时使用
\(D_0\ge-\delta_{{\rm fr},1}\)。
最大剩余增量在被放大的主范数中为
\(2\sigma-\ell_*=5\sigma/3\)；
其他范数的增量为 \(\sigma\)。因此对角
仅需要所显示的末端损失 \(5\sigma/3\)，加上
已保留的频率扰动 \(\delta_{{\rm fr},1}\)
与数值支集修正 \(\theta_N\)，当 \(A\le M\) 时，
以及在带填充的零槽核心中的一个附加 \(\delta\)。
\(\theta_N\) 修正包含在 \(C_*\xi\) 阶段余量中。

对于 \(j\ne0\)，首先在这个真正的完整
\((u,v,j)\) 和上进行整二进局部化，然后才抽取任何公共支集。
名义尺度为
\(T_2=Z^{2a_0-K-g}\)，核的幅角在其
当前 Gauss 扇区中的上确界至少为
\(q_j/(e^{2H_N}T_{\rm sec})\)。上述尾部估计适用，使用
\(|F(u,v;j)|\le q_uq_v\)，且字面条件
\(s\mid u,v\) 仍然存在。记
\[
 \Omega_{\rm ret}(q_j)
 :=\sum_{\lambda\ {\rm retained}}\omega_\lambda(q_j/T_\lambda).
\]
此公共行权重满足 \(0\le\Omega_{\rm ret}\le1\)。
方程~\eqref{old-eq:2.18b} 中的完整核保留在
每个保留项上。所有后续带符号 Fourier 分离每次
对一个保留的 \(\lambda\) 进行，因此正规化行变量
停留在一个固定对数盒上；\(\Omega_{\rm ret}\) 记录它们的和与
公共定义域。在有界多项式范围内有 \(O(\log Z)\) 个
相关的保留二进块。前面的 \(j=0\) 项正是
同一最终 \(\Phi_2\) 球的对角，包括先前加入的
Gauss 行零；两个零项都不在不同定义域上恢复。

现在记 \(u=D_2a,\ v=E_2b\)，抽取真正的完备
公共支集，使得 \((a,b)=1\) 且 \((ab,D_2E_2)=1\)。
按第一变换那样将所有抽取的幂分配给素朴变量与槽。
引理~\ref{lem:complete-support-correlation} 在此真正
轨迹上适用，给出
\[
 F(D_2a,E_2b;j)=F(D_2,E_2;j)\mathcal R(a,E_2)
   \overline{\mathcal R(b,D_2)}\mathcal R(a,b)
   \chi_a(j)\overline{\chi_b(-j)}.
\]
它允许所有四个模中的任意素数幂，并在公共素数处
留下一个行标量。每个被抽取的素数在其一侧
穿孔所有剩余因子，且供给该素数的每个槽都被
冻结。

这也显式解释了槽与素朴变量
共享素数的情形。若 \(p\) 只出现在一侧，具有
全重数 \(i\)，则其剩余因子为
\(\chi_p(j)^i\)。将 \(i\) 在两个素朴变量与
可能的槽之间拆分，由乘性恰好给出相同的因子，
包括当 \(6\mid i\) 且 \(p\mid j\) 时。例如，一个素朴
\(p^5\) 与一个槽 \(p\) 给出零扩张因子
\(\chi_p(j)^6\)。若 \(p\) 出现在两侧，则其全部幂被
移除。在不相等情形 \(i>j_0\ge1\)，局部相关性
除非 \(6\mid j_0\) 且频率为 \(p^{j_0}k\)
（其中 \(p\nmid k\)）外为零；当它非零时等于
\[
 P^{j_0-1}(P-1)\chi_p(k)^{i-j_0},\qquad P=q_p.
\]
这是一个行标量，即使第一侧的超额赋值
部分来自活跃槽。例如 \(i=7,j_0=6\) 与一个
素朴 \(p^6\) 及槽 \(p\) 留下标量
\(P^5(P-1)\chi_p(k)\)，冻结该槽，并在 \(p\) 处穿孔两个
剩余的素朴变量。相等重数具有
方程~\eqref{eq:correlation-local} 中的标量因子与
相同的穿孔结论。完备而非仅 gcd 的抽取
正是使这些局部因子独立于剩余
变量的原因。

定义
\[
 c_2=\log_Zq_{D_2},\quad d_2=\log_Zq_{E_2},\quad
 b_2=(c_2+d_2)/2.
\]
令 \(p_2\) 为其公共根式的对数范数，
令 \(G_c=(D_2,E_2)\)，并令 \(g_2=\log_Zq_{G_c}\)。
除非 \(G_c\mid j\)，相关性为零。在具有
相等重数 \(i\not\equiv0\pmod6\) 的公共
素数处，根据该素数不整除或整除相除后的频率
\(j/G_c\)，称其为单位或非单位。令 \(t_2\) 为这些
单位素数的总根式长度。令 \(V_{\rm id}\) 为
这些非单位素数的乘积并令
\(V=\log_Zq_{V_{\rm id}}\)。

方程~\eqref{eq:correlation-local} 给出以下绝对
局部界：
\[
 \begin{array}{c|c}
 \text{公共重数与相除后的频率}&
       \text{绝对相关性界}\\ \hline
 i=i,\ 6\nmid i,\ \text{unit}&P^{i-1}\\
 i=i,\ 6\nmid i,\ \text{nonunit}&P^i\\
 i=i,\ 6\mid i,\ \text{either}&P^i\\
 i>j_0,\ 6\mid j_0,\ \text{unit}&P^{j_0}.
 \end{array}
\]
不在最后一行的每个不相等情形均为零。固定
所指明的单位/非单位划分，并记 \(j=G_cV_{\rm id}h'\)。
划分标量
\[
 1_{\rm part}(h')
 \frac{F(D_2,E_2;G_cV_{\rm id}h')}{Z^{g_2-t_2}}
\]
定义为在该划分之外为零，且在此定义下
处处模至多为一。我们不断言其他行上未划分
相关性的单位界。被旧掩码强制为零的每个冻结分配
在使用此绝对界之前仍被丢弃。

完备支集恒等式在每个整剩余列上
留下行因子
\(\chi_n(G_cV_{\rm id}h')\)。
为计数其固定移动支集，令 \(v=G_cV_{\rm id}\) 并
\(e_p=v_p(v)\bmod 6\)。这些好理想的标准主
生成元的乘积是主的，因此精确的全输入
恒等式为
\[
 \chi_n(v)=
   \prod_{\substack{p\mid v\\e_p\ne0}}\chi_n(p)^{e_p}
   \prod_{\substack{p\mid v\\e_p=0}}1_{(n,p)=1},
\]
而原行的任何单位保留在 \(h'\) 中。此恒等式中的所有素数
已经穿孔每个剩余因子，因为它们属于
真正的抽取支集。活跃的 \(e_p\ne0\) 素数
包含在由 \(t_2\) 计数的单位集与由 \(V\)
计数的非单位集中。相等六可除与不相等最小六可除
素数具有 \(e_p=0\)，且在此固定因子中仅由公共
穿孔表示。非单位素数也可以具有
\(e_p=0\)，此时在 \(V\) 中计数它只会扩大
界。这只分解 \(\chi_n(v)\)：自然行因子
\(\chi_n(h')\) 及其全部零点，以及所有固定射线互反
相位均保留。它不会将随变化行的抵消重新归类
为额外穿孔。因此完整的新移动支集
长度至多为 \(\widetilde q+w_o+t_2+V\)。

定义名义行宽度与总宽度为
\begin{equation}
 \begin{aligned}
 m'&=2a_0-K-g-g_2-V,&
 q'&=\widetilde q+w_o+t_2+V,\\
 M'&=m'+q'
    =M+J-g-g_2+t_2
    \le M-\sigma.
 \end{aligned}
 \label{old-eq:2.13}
\end{equation}
这里 \(g_2\ge t_2\)，且 \(g\) 的定义给出
\(g-J\ge\sigma\)。这些恒等式通过代入
\(a_0=A-c-w\) 以及方程~\eqref{old-eq:2.3} 之前的
\(K_0\) 的定义得到。令
\[
 0\le\delta_{{\rm fr},2}:=\frac\xi2+
       \frac{\log(C_dC_{\rm sec})}{\log Z}<\xi.
\]
这里的扇区常数是第二变换的那些常数。
保留的行权重恰好拉回到
\(\Omega_{\rm ret}(q_{G_cV_{\rm id}}q_{h'})\)。在其支集上，
\[
 q_{h'}\le
 \frac{C_dT_{\rm sec}Z^{\xi/2}}{q_{G_cV_{\rm id}}}
 \le\frac{C_dC_{\rm sec}Z^{2a_0-K-g+\xi/2}}
          {q_{G_cV_{\rm id}}}
 =Z^{m'+\delta_{{\rm fr},2}}.
\]
因此它位于那个公共包络行球中。在
非空保留范围上定义所声明的非负行长度
\(m'_{\rm act}=\max\{0,m'+\delta_{{\rm fr},2}\}\) 与
\(M'_{\rm act}=m'_{\rm act}+q'\)。非空性蕴含
\(m'+\delta_{{\rm fr},2}\ge0\)，故
\(0\le M'_{\rm act}-M'\le\delta_{{\rm fr},2}\)。
在整个带符号计算中，将此包络
球其余部分的拉回权重视为零，并对
划分标量同样处理。这是一个精确的零扩张。它们
只有在形成非负子范数或例外绝对乘积之后
才会被其绝对界替换。
完备抽取仅缩短一个素朴变量或冻结
整个槽，因此存活的槽长度满足 \(z'\le z\)。

我们记录第二变换中所有其他指数。
条件 \(s\mid u,v\) 未被丢弃：它蕴含 \(s\) 的
每个素数都属于第二个完备公共根式。因此
完备抽取之后，剩余列上不存在剩余的
\(s\) 整除条件。对于此固定的
\(s\)，长度为 \(p_2\) 的可能根式个数为 \(O(Z^{p_2-s_0+\epsilon_1})\)：将该根式写作
\(s\mathfrak r_2\)，其中 \((s,\mathfrak r_2)=1\)，并计数
剩余范数的无平方因子 \(\mathfrak r_2\)。若在固定壳边界之外
\(p_2<s_0\)，则没有。用于 \(C,D\) 的同一 Rankin
论证界定其幂与分配重数。选择单位/非单位
划分的代价至多为
\(2^{\omega(\operatorname{rad}(D_2E_2))}\)，这是另一个除子有界
因子，包含在 \(\epsilon_1\) 中。这些指数为
\[
 \begin{array}{l|r}
 \text{factor}&\text{exponent}\\ \hline
 \text{平方中心正规化}&-a_0\\
 \text{行 Poisson 因子与正规化逆根}&K+g-a_0\\
 \text{素密度（当出现时）}&-\ell\\
 \text{固定 }s\text{ 的公共支集计数}&p_2-s_0\\
 \text{common correlation}&g_2-t_2\\
 \text{转换为正规化剩余乘积}&a_0-b_2
 \end{array}
\]
最后一行为
\(\{(a_0-c_2)+(a_0-d_2)\}/2\)。
它们的和为
\[
 K+g-\ell-a_0+p_2-s_0+g_2-t_2-b_2.
\]
从余量
方程~\eqref{old-eq:2.7} 中减去此和，得到
内层素朴乘积所允许的指数。这是要由
更小宽度矩或由例外行估计提供的界：
\begin{equation}
 M'+\Delta_{\rm child},\qquad
 \Delta_{\rm child}=b_2-p_2+w+\mathcal B_c+\ell
                   \ge w+\ell.
 \label{old-eq:2.14}
\end{equation}
该不等式使用 \(b_2\ge p_2\)，因为每个公共素数处的
两个重数都至少为一。

我们现在在形成独立子之前移除剩余互质性。
对于固定的真正 \(D_2,E_2\)，定义
\[
 \widetilde F_{D_2,E_2}(a,b;j)
 :=F(D_2,E_2;j)\mathcal R(a,E_2)
   \overline{\mathcal R(b,D_2)}\mathcal R(a,b)
   \chi_a(j)\overline{\chi_b(-j)}
\]
在所有各自与 \(D_2E_2\) 互素且被旧掩码允许的
剩余对上。仅当 \((a,b)=1\) 时它等于
真正的相关性。
通过所分配的卷积公式与旧公共掩码，在所有
各自允许的剩余对上定义剩余的
\(\mathcal C_2(a,b;j)\)，包括 \((ab,D_2E_2)=1\)。
保持形式整乘积范数
\(q_{D_2}q_a,q_{E_2}q_b\)、两个逆根、完整光滑
核，以及在其包络球上以零扩张的相同拉回行权重。
对每个保留行，精确恒等式为
\[
 \begin{split}
 \sum_{(a,b)=1}\mathcal C_2(a,b;j)F(D_2a,E_2b;j)
 &=\sum_{\mathfrak t\ {\rm squarefree}}\mu(\mathfrak t)
   \sum_{\mathfrak t\mid a,b}\mathcal C_2(a,b;j)
                \widetilde F_{D_2,E_2}(a,b;j).
 \end{split}
\]
这是完整的 M\"obius 指示函数，而非截断。各个
列壳与保留行球是有限的。因此其相等性
通过交换有限和并仅在 \((a,b)=1\) 时
使用真正相关性恒等式得到。\(\mathfrak t\) 的额外公共
素数不属于真正 \(D_2,E_2\) 支集或
其频率限制。由于剩余掩码已经
排除该支集，每个非零 \(\mathfrak t\) 与它及与
\(s\) 都不相交。

在每个保留二进块上，如第一变换那样，在
正规化 \(h'\) 范数与两个整乘积范数中分离完整核与逆根。
冻结因子 \(G_cV_{\rm id}\) 只对行 Fourier 幂贡献
一个外部相位。固定支集盒在当前活跃标签之前
选定，因此所得到的系数测度对它们、对
两个矩形以及对行都是公共的。然后固定 \(\mathfrak t\)。
对其每个素数 \(p\)，
使用精确的因子分配恒等式
\[
 1_{p\mid\prod_i n_i}
 =\sum_{\varnothing\ne J}(-1)^{|J|+1}
       \prod_{i\in J}1_{p\mid n_i},
\]
其中因子 \(n_i\) 是两个素朴变量与活跃
槽。对选中的素朴变量记 \(n_i=pl_i\)，对 \(l_i\) 无
限制。这只抽取一个行标量及
\(q_p^{it}\)，并将两个 \(X_i,Y_i\) 替换为
\(X_i/q_p,Y_i/q_p\)。它不引入单变量穿孔。
选中的槽被冻结。若两个不同的除子素数选中
同一个素数槽，则该为零。商与未选中的
因子仍可包含 \(p\)，并可与对侧重叠；
不施加新的互质性。旧公共掩码保持公共，且
两个新素朴尺度的乘积在两个矩形中相等。
抽取的行标量，包括零标量，保留在
带符号恒等式中。

只有在现在，对于固定的保留二进块与固定 Fourier 参数，
每个分离出的被加项才是两个独立剩余
卷积的乘积，其中 \(\mathbf J\) 记录因子分配。
它们的行特征为
\[
 \tau_1(n)\rho(n)\chi_n(G_cV_{\rm id}h'),\qquad
 \tau_1(n)\rho'(n)\chi_n(-G_cV_{\rm id}h'),
 \quad \rho,\rho'\in\Theta.
\]
它们通过上述固定因子支集表示来理解。特别地
它们的诱导特征之比属于 \(\Theta\)，包括补充因子
\(n\mapsto\chi_n(-1)\)；此断言在零点处不取商。
它们的自然行零点被保留。抽取
选中的因子贡献行标量，而非剩余特征的
新因子，也不是 \(q'\) 中的新支集。这些标量，
包括零点，保留在带符号恒等式中。

\begin{lemma}[完备抽取下的公共系数]
\label{lem:centered-coefficient-invariant}
对于中心化输入，两个变换以及刚构造的中间 Gauss 行
放大保持以下系数数据。
每个 Gauss 多项式与每个放大器误差保持
方程~\eqref{eq:centered-coefficient-form}，乘以 \(G(u,h)\)，
除有界冻结标量与模至多为一的行标量外。
对于非中心化输入，全程省略第二个矩形。
在第二变换与完整 M\"obius 因子分配之后，每个
分离的子侧具有该系数形式，但不含 \(G\) 且不含
\(s\mid u\)。每个存活槽具有子侧的一个整乘积
特征及其原始系数 \(\nu_i\)，不含 Gauss 系数。
在一侧内，素朴变量在两个矩形中具有相同的特征、穿孔掩码
与范数幂。一个平方 Gauss 范数的两侧
可以有不同的范数幂；它们的诱导特征相差
\(\Theta\) 中的一个成员。

若新的诱导特征属于 \(\Theta\)，则固定活跃槽
标签留下素朴系数
\[
 \vartheta(l_1)\vartheta(l_2)
 1_{(l_1l_2,\mathfrak R_*)=1}q_{l_1l_2}^{it}D_{\mathbf b}(l_1,l_2)
\]
对某个 \(\vartheta\in\Theta\)、某个多项式范数的无平方因子掩码
\(\mathfrak R_*\) 及某个实数 \(t\)。
掩码可以依赖于冻结行，但对两个变量
与两个矩形都是公共的。所有系数都以其保留的
冻结标量（包括零点）以及下面记录的精确正规化来理解。
\end{lemma}

\begin{proof}
第一变换的计算证明了 Gauss 系数形式；局部
\(p\) 幂计算对每个放大器误差证明了它。
上述真正的完备支集恒等式与完整 M\"obius 分配
证明了子形式，在每个整剩余乘积上
保持一个特征与范数幂。因此“公共”涉及一个分离的侧；
所显示的特征比是两侧之间的关系。
当诱导特征属于 \(\Theta\) 时，它在单位上等于某个
\(\vartheta\in\Theta\)。其冗余的自然零点
与固定穿孔合并为一个多项式范数的
\(\mathfrak R_*\)。对每个固定的 \(\Pi\)，乘性把它在
\(\Pi l_1l_2\) 上的值、掩码与范数幂分解为 \(\Pi\)
中的一个标量乘以结论中陈述的三个因子。这证明了
例外断言。
\end{proof}
